\hspace{3mm}

\centerline{\bf \Huge 6 класс}


\problem{\textbf{1}} 
В ЭлМШ учатся Шиншиллы, которые всегда лгут, Лососи, которые всегда говорят правду, и Киви, которые могут сказать что угодно по своему усмотрению. Однажды 20262026 учеников ЭлМШ построились \textit{черепахой} (прямоугольник из 10001 строки и 2026 столбцов). Затем каждый сказал: «В строке, в которой я стою, Шиншилл больше, чем остальных вместе взятых в этой строке, а в столбце меньше, чем остальных вместе взятых в этом столбце». Какое наименьшее и наибольшее количество Шиншилл может быть в данном построении \textit{черепахой}?

\problem{\textbf{2}} 
Расул Владимирович придумал новую операцию над числом под названием «\textit{СилаЭМГ}». Операция «\textit{СилаЭМГ}»: I. Если число однозначное, то заменяет его на число в 6 раз большее; II. Если число не однозначное, то заменяет его на число в 6 раз большее и после стирает последнюю цифру этого числа. Какое число получит Анджур Аркадьевич после 2025 применений операций «\textit{СилаЭМГ}» над числом 2026.

\problem{\textbf{3}} 
Марат загадал два различных натуральных числа, произведение которых больше 14. Лера загадала два различных натуральных числа, больших чем числа Марата. Данзан загадал ещё два различных натуральных числа, больших чем числа Леры, и произведение которых меньше 89. Найдите все возможные варианты этих шести чисел и докажите, что других вариантов быть не может.

\problem{\textbf{4}} 
Вдоль «\textit{сладкой}» стены в ЭТЛ стоят 12 сундуков, внутри каждого из которых по 600 чокопаек. Разрешено забирать из сундука ровно 100, ровно 200 или ровно 300 чокопаек, но при этом никто не должен забирать поровну чокопаек с соседних сундуков, иначе Анджур Аркадьевич отчислит из ЭТЛ. София, потом Дима и после них Энкира прошли вдоль «\textit{сладкой}» стены и забрали абсолютно все чокопайки. Дима забрал меньше 1900 чокопаек. Докажите, что одна из двух других девочек забрала не менее 3000 чокопаек.

\problem{\textbf{5}} 
4 января 2026 года 52 ученика ЭлМШ находились в Калмыкии. С помощью двухместного \textit{Мусорного пакета} (вид транспорта, похожий на двухместный автомобиль) все ученики переместились в Дагестан (некоторые рейсы, возможно, выполнялись в одиночку). 
Могло ли случиться, что каждая пара учеников была в пути ровно один раз (с Калмыкии в Дагестан или наоборот)?

\textit{\textbf{p.s.}} У \textit{Мусорного пакета} всегда должен быть водитель!
